\begin{lemma}[Shear descent in a fixed weight]
\label{lem:shear-descent}
Let \(V=A\oplus B\), with \(B\ne0\), and let \(m\ge1\).
Suppose that \(N\subset D\otimes\Sym^mV\) is invariant under
\(\mathrm{id}_A\oplus t\,\mathrm{id}_B\), \(t\in\C^*\),
and under all shears
\[
 u_L(a+b)=a+b+L(a),\qquad L\in\Hom(A,B).
\]
If \(N\not\subset D\otimes\Sym^mB\), then
\[
 N\cap\bigl(D\otimes A\otimes\Sym^{m-1}B\bigr)\ne0.
\]
More precisely, a nonzero vector in the \(A\)-degree-\(k\)
weight space, \(k>0\), has a nonzero image in the degree-one
space under a product of \(k-1\) infinitesimal rank-one shears.
The same fixed nonzero vector of \(B\) can be used in every shear.
\end{lemma}
\begin{proof}
Put \(s=\dim B\).  The displayed torus acts on
\(D\otimes\Sym^kA\otimes\Sym^{m-k}B\) with weight
\(s+m-k\).  These weights are distinct.  Since a complex torus
acts semisimply, its invariant subspace \(N\) is the direct sum
of its intersections with these weight spaces.  This justifies
projecting a vector of \(N\) to one degree before differentiating;
no cancellation between degrees is possible.

Choose a basis \(a_1,\ldots,a_r\) of \(A\), a nonzero
\(b\in B\), and a basis vector \(\epsilon\) of \(D\).
Write a nonzero degree-\(k\) vector as \(\epsilon\otimes F\),
where
\[
 F=\sum_{|\beta|=k}a^\beta f_\beta,\qquad
                 f_\beta\in\Sym^{m-k}B.
\]
Choose \(\beta\) with \(f_\beta\ne0\), and \(i\) with
\(\beta_i>0\).  Set \(\alpha=\beta-e_i\).  In
\(\partial_A^\alpha F\), the coefficient of \(a_i\) is
exactly \(\beta!f_\beta\).  Indeed a monomial that contributes
to this coefficient must have exponent \(\alpha+e_i=\beta\).
This coefficient is nonzero in characteristic zero, so the derivative
is a nonzero element of \(A\otimes\Sym^{m-k}B\).

The shear sending \(a_j\) to \(a_j+tb\) and fixing the other
basis vectors has determinant one.  Its infinitesimal action on
\(D\otimes\Sym^mV\) is consequently
\(1_D\otimes b\partial_{a_j}\), with no extra action on
\(D\).  Invariance under the shear implies invariance under its
infinitesimal action: the orbit is a polynomial in \(t\), and
all its coefficients belong to the invariant linear subspace.
These differential operators commute, because every
\(\partial_{a_j}\) kills \(b\).  Their product with
multiplicities \(\alpha_j\) sends our vector to
\[
 \epsilon\otimes b^{k-1}\partial_A^\alpha F.
\]
This vector belongs to \(N\).  It is nonzero, since multiplication
by \(b^{k-1}\) is injective in the polynomial algebra \(\Sym B\).
Its exact bidegree is \((1,m-1)\); thus its image in the corresponding
weight-graded piece is also nonzero.  The proof uses neither a special
choice of \(b\) nor a restriction on \(r\) or \(s\).  When
\(k=1\) the product of shears is the identity.  This proves all
assertions, including the low-dimensional cases.
\end{proof}

